\documentclass[10pt,a4paper]{amsart}
\usepackage[margin=19mm]{geometry}
\usepackage{amsmath,amssymb,mathtools}
\usepackage[scr=rsfs,cal=euler]{mathalfa}
\usepackage{xfrac}
\usepackage[unicode,hidelinks,bookmarks=false]{hyperref}
\newcommand{\ie}{i.e.\ }
\newcommand{\cf}{cf.\ }
\newcommand{\resp}{respectively\ }
\newcommand{\Subsection}{\subsection}
\providecommand{\nobreakdash}{\nobreak\hbox{-}}
\setlength{\emergencystretch}{2em}
\multlinegap0pt
\usepackage{xcolor}
\usepackage{amssymb}
\usepackage{tikz}
\newtheorem{proposition}{Proposition}
\newtheorem{lemma}{Lemma}
\def \R {\mathbb{R}}
\title{A conformally invariant Dirac-type equation on compact spin manifolds: the effect of the geometry.}
\author{Ali Maalaoui, Vittorio Martino}
\date{Source: arXiv:2604.08738v1 (CC BY 4.0)}
\begin{document}
\maketitle

\noindent\emph{Source and licence.} A conformally invariant Dirac-type equation on compact spin manifolds: the effect of the geometry.. Version arXiv:2604.08738v1 (CC BY 4.0), distributed by arXiv under the Creative Commons Attribution 4.0 International licence, http://creativecommons.org/licenses/by/4.0/, as recorded in the arXiv licence metadata for this exact version and shown on the abstract page https://arxiv.org/abs/2604.08738v1. Full licence notice: \texttt{licenses/arxiv-2604.08738-CC-BY-4.0.txt}.\par\smallskip

\noindent\textit{Editorial scope.} Counted unit: the article's proposition exp (source0.tex lines 597--607). The article's complete original proof of it is delivered with it (source0.tex lines 609--764, one \texttt{proof} environment of 156 lines). No mathematics is written, repaired or re-derived: the statement and the proof are the article's own lines, in the article's own notation, and no retained line is substituted or re-flowed.  Retained uncounted as the source-local context the proof needs: lines 351--362; lines 450--452; lines 456--458; lines 554--564; lines 570--580; lines 586--591.  Nothing was omitted from any retained span, so the package carries no omission record.  No retained line contains a citation, so the wrapper carries no bibliography and no bibliography entry.  No retained line contains a figure environment, an image-inclusion command or drawing code, so no image is delivered and the package declares no asset dependency.  Editorial formatting only, in addition: an AMS \texttt{amsart} preamble that restates the article's own declarations and macros used by the retained lines, namely the environments \texttt{proposition}, \texttt{lemma} on separate counters, together with the standard packages those lines need.  The retained environments print in this document as Proposition 1, Lemma 1 in order of appearance, whereas the article prints the same items with the numbering of its own full document.  The complete native source of the article as distributed in the arXiv e-print for this exact version is bundled unchanged as \texttt{sources/198167/source0.tex}.
\begin{proposition}\label{propgreen}
We consider a compact Riemannian manifold $(M,g)$ as above with positive Yamabe constant $Y(M,[g])>0$, then the Green's function $G_{g}$ of the conformal Laplacian $L_{g}$ has the following asymptotic expansion, in a conformal normal coordinate system centered at a point $p$, that we can identify with the origin, with $r=|x|$.
\begin{itemize}
\item[i)] For $n=3,4,5$ or $(M,g)$ locally conformally flat around a point $p$: % that we can identify with the origin and $r=|x|$:
$$G_{g}(p,x)=b_{n}\Big(\frac{1}{r^{n-2}}+A(p)\Big)+o(1).$$
\item[ii)] For $n=6$
$$G_{g}(p,x)=b_{n}\Big(\frac{1}{r^{4}}-\frac{1}{1440}|W(p)|^{2}\ln(r)\Big)+O(1).$$
\item[iii)] For $n\geq 7$
$$G_{g}(p,x)=\frac{b_{n}}{r^{n-2}}\Big(1+\frac{(n-2)}{48(n-1)(n-4)}\Big(\frac{r^{4}}{12(n-6)}|W(p)|^{2}-R_{,ij}x^{i}x^{j}r^{2}\Big)\Big)+O(r^{7-n}).$$
where $b_{n}=\frac{1}{(n-2)\omega_{n-1}}$, $W(p)$ is the Weyl tensor at $p$ and $R_{,ij}$ is the covariant derivative of the scalar curvature $R=R_g$ with respect to $\partial_{x_{i}}$ and $\partial_{x_{j}}$.
\end{itemize}
\end{proposition}

\begin{equation}\label{psi}
\Psi=\Big(\frac{1}{1+|x|^{2}}\Big)^{\frac{n}{2}}(1-x)\cdot \Psi_{0},
\end{equation}

\begin{equation}\label{psiepsilon}
\psi_{\varepsilon}(x)=\eta(x)\varepsilon^{-\frac{n-1}{2}}\Psi\left(\frac{x}{\varepsilon}\right)=\eta(x) \Psi_{\varepsilon}(x).
\end{equation}

\begin{proposition}\label{propest1}
In local coordinates we have
\begin{itemize}
\item[i)] $b_{ij}-\delta_{ij}=-\frac{1}{6}R_{i\alpha\beta j}x^{\alpha}x^{\beta}-\frac{1}{12}R_{i\alpha\beta j,k}x^{\alpha}x^{\beta}x^{k}+O(r^{4})$.
\item[ii)] $X=-\sum\limits_{k=1}^{n}\Big(\frac{1}{4}R_{\alpha k}x^{\alpha}+\frac{1}{6}R_{\alpha k,\beta}x^{\alpha}x^{\beta}+O(r^{3})\Big)e_{k}$.
\item[iii)] $\Omega=-\frac{1}{144}\sum\limits_{\begin{array}{cc}
i,j,k\\
 i\not=j\not=k\not=i
\end{array}}\sum\limits_{\ell}R_{\ell \beta \gamma k}(R_{ji\alpha\ell}+R_{j\ell\alpha i})x^{\alpha}x^{\beta}x^{\gamma}e_{i}\cdot e_{j}\cdot e_{k}+O(r^{4})$.
\end{itemize}
\end{proposition}

\begin{proposition}\label{proplp} {\textcolor{red}{ 22/9; tau 1/9 }}
Given $p\in M$ and $N\geq 2$. Then there exits a conformal metric $g$ on $M$ such that around the point $p$ we have
$$det(g_{ij})=1+O(r^{N}).$$
Moreover, at the point $p$, it holds
\begin{itemize}
\item[i)] $R_{ij}=0$.
\item[ii)] $R_{ij,k}+R_{jk,i}+R_{ki,j}=0$.
\item[iii)] $\Big(R_{\alpha \beta,k\ell}+\frac{2}{9}\sum \limits_{i,d}R_{i\alpha \beta d}R_{ik\ell d}\Big)x^{\alpha}x^{\beta}x^{k}x^{\ell}=0.$
\item[iv)] $-\Delta R=-R_{,kk}=\frac{1}{6}|W|^{2}$.
\end{itemize}
\end{proposition}

\begin{lemma}\label{lemyx}
For $\Psi$ defined as in $(\ref{psi})$, we have for any $x\in \R^n$
$$R_{i\alpha\beta j}x^{\alpha}x^{\beta}\partial_{x_{i}}\cdot \nabla _{\partial_{x_{j}}}\Psi (x) =0.$$
Moreover, if $A_{ijk\ell}\in \R$ with $1\leq i,j,k,\ell \leq n$, then we can choose $\Psi_{0}$ in $(\ref{psi})$, so that
$$A_{ijk\ell}\langle \partial_{x_{i}}\cdot \partial_{x_{j}}\cdot \partial_{x_{k}}\cdot \partial_{x_{\ell}}\cdot \Psi_{0},\Psi_{0}\rangle =0.$$
\end{lemma}

\begin{proposition}\label{lemma1 exp}
For $\varphi_{\varepsilon}$ defined as above, we have
$$\|\nabla J_{g}(\varphi_{\varepsilon})\|_{H^{-\frac{1}{2}}}\leq \left\{\begin{array}{ll}
O(\varepsilon^{\frac{n-1}{2}}), &  4\leq n\leq 6\\
\\
O(\varepsilon^{3}|\ln(\varepsilon)|^{\frac{4}{7}}), &  n=7\\
\\
O(\varepsilon^{3}), &  n\geq 8
\end{array}
\right. .$$
\end{proposition}

\begin{proof}
We will proceed by estimating first the $H^{-\frac{1}{2}}$-norm of $\varphi_{\varepsilon}$ and then the norm of 
$$\nabla J_{g}(\varphi_{\varepsilon})=\overline{D}\varphi_{\varepsilon}-\Big(V_{g}*|\varphi_{\varepsilon}|^{2}\Big)\varphi_{\varepsilon}\; .$$ 
Here, $H^{-\frac{1}{2}}$ is the dual of the space $H^{\frac{1}{2}}(M)$ equipped with the norm $\|\cdot \|_{\frac{1}{2}}$. By the continuous embedding $L^{\frac{2n}{n+1}}(M)\hookrightarrow H^{-\frac{1}{2}}(M)$, we have

\begin{align}
\|\varphi_{\varepsilon}\|_{H^{-\frac{1}{2}}}&\leq C \|\varphi_{\varepsilon}\|_{L^{\frac{2n}{n+1}}}=\Big(\int_{B_{2\delta}}|\varphi_{\varepsilon}|^{\frac{2n}{n+1}}\ dv\Big)^{\frac{n+1}{2n}}\notag\\
&\leq C \Big(\int_{|x|\leq 2\delta} |\psi_{\varepsilon}(x)|^{\frac{2n}{n+1}}\ dx \Big)^{\frac{n+1}{2n}}\notag\\
&\leq C\varepsilon\Big(\int_{0}^{\frac{2\delta}{\varepsilon}} \frac{r^{n-1}}{(1+r^{2})^{\frac{n(n-1)}{n+1}}}\ dr \Big)^{\frac{n+1}{2n}} \notag\\
&\leq C\left\{\begin{array}{ll}
\varepsilon |\ln(\varepsilon)|^{\frac{2}{3}}, & n=3\\
\\
\varepsilon, & n\geq 4 
\end{array}
\right..\label{phieps}
\end{align}
Next, we focus on estimating $\nabla J_{g}(\varphi_{\varepsilon})$. Indeed, we have
\begin{align}
\overline{D}\varphi_{\varepsilon}&=\overline{D\psi_{\varepsilon}}+\Omega\cdot \overline{\psi_{\varepsilon}}+X\cdot \overline{\psi}+\sum_{i,j}(b_{ij}-\delta_{ij})\overline{\partial_{x_{i}}\cdot \nabla_{\partial_{x_{j}}}\psi_{\varepsilon}}\notag\\
&=\Big(\int_{\R^{n}}V_{g_{\R^{n}}}(x,y)|\Psi_{\varepsilon}(y)|^{2}\ dy\Big) \varphi_{\varepsilon}+(\nabla\eta(x)+X) \cdot \varphi_{\varepsilon}+\Omega\cdot \varphi_{\varepsilon}\notag\\
&\quad+\sum_{i,j}(b_{ij}-\delta_{ij})\overline{\partial_{x_{i}}\cdot \nabla_{\partial_{x_{j}}}\psi_{\varepsilon}}.\label{decomp}
\end{align}


\noindent
On the other hand,
\begin{align}
\overline{D \psi_{\varepsilon}}-(V_{g}*|\varphi_{\varepsilon}|^{2})\varphi_{\varepsilon}&=\Big(\int_{\R^{n}}V_{g_{\R^{n}}}(x,y)|\Psi_{\varepsilon}(y)|^{2}\ dy \Big)\varphi_{\varepsilon} \notag\\ 
&\quad - \Big(\int_{M}V_{g}(x,y)|\varphi_{\varepsilon}|^{2}\ dv(y)\Big) \varphi_{\varepsilon}+\nabla\eta(x) \cdot \varphi_{\varepsilon}\notag\\
&=\Big(\int_{|x-y|<\frac{\delta}{2}}[V_{g_{\R^{n}}}(x,y)-V_{g}(x,y)]|\varphi_{\varepsilon}(y)|^{2}\ dy \Big) \varphi_{\varepsilon}\notag\\
&\quad + \Big(\int_{|x-y|<\frac{\delta}{2}}V_{g}(x,y)O(|y|^{N})|\varphi_{\varepsilon}(y)|^{2}dy\Big)\varphi_{\varepsilon}\notag\\
&\quad + \Big(\int_{|x-y|>\frac{\delta}{2}}V_{g_{\R^{n}}}(x,y)|\Psi_{\varepsilon}(y)|^{2} dy\Big) \varphi_{\varepsilon}\notag\\
&\quad - \Big(\int_{|x-y|>\frac{\delta}{2}}V_{g}(x,y)|\varphi_{\varepsilon}|^{2}\ dv(y) \Big) \varphi_{\varepsilon} +\nabla \eta \cdot \varphi_{\varepsilon}.\notag
\end{align}



\noindent
This leads to 
\begin{align}
\overline{D}\varphi_{\varepsilon}-(V_{g}*|\varphi_{\varepsilon}|^{2})\varphi_{\varepsilon}&=\Big(\int_{|x-y|<\frac{\delta}{2}}[V_{g_{\R^{n}}}(x,y)-V_{g}(x,y)]|\varphi_{\varepsilon}(y)|^{2}\ dy \Big) \varphi_{\varepsilon}\notag\\
&\quad+\Big(\int_{|x-y|>\frac{\delta}{2}}V_{g_{\R^{n}}}(x,y)|\Psi_{\varepsilon}(y)|^{2} dy\Big) \varphi_{\varepsilon}-\Big(\int_{|x-y|>\frac{\delta}{2}}V_{g}(x,y)|\varphi_{\varepsilon}|^{2}\ dv(y) \Big) \varphi_{\varepsilon} \notag\\
&\quad+\Big(\int_{|x-y|<\frac{\delta}{2}}V_{g}(x,y)O(|x-y|^{N})|\varphi_{\varepsilon}(y)|^{2}dy\Big)\varphi_{\varepsilon}\notag\\
&\quad +\nabla \eta \cdot \varphi_{\varepsilon} + \Omega\cdot \overline{\psi_{\varepsilon}}+X\cdot \overline{\psi}+\sum_{i,j}(b_{ij}-\delta_{ij})\overline{\partial_{x_{i}}\cdot \nabla_{\partial_{x_{j}}}\psi_{\varepsilon}}\notag\\
&=A_{1}+A_{2}+A_{3}+A_{4}+A_{5}+A_{6}+A_{7}+A_{8}.\notag
\end{align}
We will estimate now the terms $A_{i}$ for $i=1,\cdots,8$, in dimensions $n\geq 4$. Indeed, for $A_{5}$, we have
\begin{align}
\|A_{5}\|_{H^{-\frac{1}{2}}}&\leq C\|A_{5}\|_{L^{\frac{2n}{n+1}}}=C\Big(\int_{B_{2\delta}}|\overline{\nabla \eta \cdot \Psi_{\varepsilon}}|^{\frac{2n}{n+1}}\ dv \Big)^{\frac{n+1}{2n}}\notag\\
&\leq C\Big(\int_{\delta\leq |x|\leq 2\delta}|\Psi_{\varepsilon}(x)|^{\frac{2n}{n+1}}\ dx \Big)^{\frac{n+1}{2n}}\notag\\
&\leq C \varepsilon \Big(\int_{\frac{\delta}{\varepsilon}}^{\frac{2\delta}{\varepsilon}}\frac{r^{n-1}}{(1+r^{2})^{\frac{n(n-1)}{(n+1)}}}\ dr \Big)^{\frac{n+1}{2n}}\leq C \varepsilon^{\frac{n-1}{2}}.\notag
\end{align}
For $A_{1}$, we use Proposition \ref{propgreen}  in Section 2, in a small coordinate patch, we have 
$$|V_{g}(x,y)-V_{g_{\R^{n}}}(x,y)|=|h(x,y)|,$$ 
with $h$ having different behavior depending on the dimension $n$. Indeed,
\begin{itemize}
\item If $n\geq 7$ then $|h(x,y)|\leq Cr^{2}$. In this case, we have
\end{itemize}
\begin{align}
\|A_{1}\|_{H^{-\frac{1}{2}}}&\leq C\|A_{1}\|_{L^{\frac{2n}{n+1}}}\leq C\Big(\int_{B_{2\delta}}\Big(\int_{B_{2\delta}}h(x,y)|\Psi_{\varepsilon}(y)|^{2}\ dy |\Psi_{\varepsilon}(x)|\ \Big)^{\frac{2n}{n+1}}dx\Big)^{\frac{n+1}{2n}}\notag\\
&\leq C\Big(\int_{B_{2\delta}}\Big(\int_{B_{2\delta}}|y|^{2}|\Psi_{\varepsilon}(x-y)|^{2}\ dy |\Psi_{\varepsilon}(x)|\Big)^{\frac{2n}{n+1}}\ dx\Big)^{\frac{n+1}{2n}} \notag\\
&\leq  C \varepsilon^{3}.\notag
\end{align}
\begin{itemize}
\item If $n=6$, then we have $|h(x,y)|\leq Cr^{2}|\ln(r)|$. Thus,
\end{itemize}
\begin{align}
\|A_{1}\|_{H^{-\frac{1}{2}}}&\leq C\|A_{1}\|_{L^{\frac{2n}{n+1}}}\leq C\Big(\int_{B_{2\delta}}\Big(\int_{B_{2\delta}}|y|^{2}|\ln(|y|)||\Psi_{\varepsilon}(x-y)|^{2}\ dy |\Psi_{\varepsilon}(x)|\Big)^{\frac{2n}{n+1}}\ dx\Big)^{\frac{n+1}{2n}} \notag\\
&\leq  C \varepsilon^{3}|\ln(\varepsilon)|.\notag
\end{align}

\begin{itemize}
\item Finally, for $n=4$ or $5$, then $|h(x,y)|\leq C$. Therefore, a similar process as above yields 
\end{itemize}
$$\|A_{1}\|_{H^{-\frac{1}{2}}}\leq C \varepsilon^{2}.$$



\noindent
Notice that a similar inequality holds for $\|A_{4}\|_{H^{-\frac{1}{2}}}$, if we choose $N\geq 4$. Hence we move to $A_{2}$. We use the fact that the Green's function is bounded outside of the diagonal and we have
\begin{align}
\|A_{2}\|_{H^{-\frac{1}{2}}}&\leq C \Big(\int_{|x|<\frac{\delta}{4}}\Big(\int_{|x-y|>\frac{\delta}{2}}|\Psi_{\varepsilon}(y)|^{2}\ dy |\Psi_{\varepsilon}(x)|^{2}\Big)^{\frac{2n}{n+1}}\ dx \Big)^{\frac{n+1}{2n}}\notag\\
&+\Big(\int_{\frac{\delta}{4}<|x|<2\delta}\Big(\int_{|x-y|>\frac{\delta}{2}}|\Psi_{\varepsilon}(y)|^{2}\ dy |\Psi_{\varepsilon}(x)|^{2}\Big)^{\frac{2n}{n+1}}\ dx \Big)^{\frac{n+1}{2n}}\notag\\
&\leq C \varepsilon\Big( \|\varphi_{\varepsilon}\|_{L^{\frac{2n}{n+1}}}\int_{\frac{\delta}{4\varepsilon}}^{\infty}\frac{r^{n-1}}{(1+r^{2})^{n-1}}\ dr +\|\varphi_{\varepsilon}\|_{L^{2}}^{2}\Big(\int_{\frac{\delta}{4\varepsilon}}^{\frac{2\delta}{\varepsilon}} \frac{r^{n-1}}{(1+r^{2})^{\frac{n(n-1)}{n+1}}}\ dr \Big)^{\frac{n+1}{2n}}\Big)\notag\\
& \leq C \varepsilon^{\frac{n+1}{2}}.\notag
\end{align}
A similar inequality holds for $\|A_{3}\|_{H^{-\frac{1}{2}}}$. Now, using Proposition \ref{propest1}, we have $|\Omega|=O(|x|^{3})$. Therefore,
\begin{align}
\|A_{6}\|_{H^{-\frac{1}{2}}}&\leq C \Big(\int_{B_{2\delta}}|\Omega|^{\frac{2n}{n+1}}|\varphi_{\varepsilon}|^{\frac{2n}{n+1}}\ dv \Big)^{\frac{n+1}{n}} \leq C \Big( \int_{|x|\leq 2\delta} |x|^{\frac{6n}{n+1}}|\Psi_{\varepsilon}(x)|^{\frac{2n}{n+1}}\ dx\Big)^{\frac{n+1}{2n}} \notag\\
&\leq C \varepsilon^{4} \Big( \int_{0}^{\frac{2\delta}{\varepsilon}} \frac{r^{\frac{6n}{n+1}+n-1}}{(1+r^{2})^{\frac{n(n-1)}{n+1}}}\ dr \Big)^{\frac{n+1}{2n}}\notag\\
&\leq C \left\{\begin{array}{ll}
\varepsilon^{\frac{n-1}{2}}, & 3\leq n\leq 8\\
\\
\varepsilon^{4}|\ln(\varepsilon)|^{\frac{5}{9}}, & n=9\\
\\
\varepsilon^{4}, & n\geq 10
\end{array}
\right..\notag
\end{align}
For $A_{7}$, working again in conformal normal coordinates and using Proposition \ref{propest1} and \ref{proplp}, we have, $|X|=O(|x|^{2})$. Hence,
\begin{align}
\|A_{7}\|_{H^{-\frac{1}{2}}}&\leq C \Big(\int_{B_{2\delta}}|X|^{\frac{2n}{n+1}}|\varphi_{\varepsilon}|^{\frac{2n}{n+1}}\ dv \Big)^{\frac{n+1}{2n}}\notag\\
&\leq C \Big(\int_{|x|\leq 2\delta}|x|^{\frac{4n}{n+1}}|\psi_{\varepsilon}(x)|^{\frac{2n}{n+1}}\ dx \Big)^{\frac{n+1}{2n}}\notag\\
&\leq C  \varepsilon^{3} \Big(\int_{0}^{\frac{2\delta}{\varepsilon}} \frac{r^{\frac{4n}{n+1}+n-1}}{(1+r^{2})^{\frac{n(n-1)}{n+1}}}\ dx \Big)^{\frac{n+1}{2n}}\notag\\
&\leq C \left\{\begin{array}{ll}
\varepsilon^{\frac{n-1}{2}}, & 4\leq n \leq 6\\
\\
\varepsilon^{3}|\ln(\varepsilon)|^{\frac{4}{7}}, & n=7\\
\\
\varepsilon^{3}, & n\geq 8
\end{array}
\right..\notag
\end{align}
It remains now to estimate $A_{8}$. Recalling that
$$A_{8}=\sum_{i,j}(b_{ij}-\delta_{ij})\overline{\partial_{x_{i}}\cdot \nabla_{\partial_{x_{j}}}\psi_{\varepsilon}} \; ,$$
from (\ref{psiepsilon}), we will write $A_{8}=B_{1}+B_{2}$, where 
$$B_{1}:=\eta\sum_{i,j}(b_{ij}-\delta_{ij})\overline{\partial_{x_{i}}\cdot \nabla_{\partial_{x_{j}}}\Psi_{\varepsilon}}\quad \text{ and  }\quad B_{2}:=\sum_{i,j}(b_{ij}-\delta_{ij})(\partial_{x_{j}}\eta)\overline{\partial_{x_{i}}\cdot \Psi_{\varepsilon}}.$$
Notice that since $|\nabla \Psi(r)|\leq C \big(f(r)\big)^{\frac{n}{2}}$, and using Lemma \ref{lemyx}, we have
\begin{align}
\|B_{1}\|_{H^{-\frac{1}{2}}}&\leq C \Big(\int_{|x|\leq 2\delta}|x|^{\frac{6n}{n+1}}|\nabla \Psi_{\varepsilon}(x)|^{\frac{2n}{n+1}}\ dx \Big)^{\frac{n+1}{2n}}\notag\\
& \leq C \varepsilon^{3}\Big(\int_{0}^{\frac{2\delta}{\varepsilon}}\frac{r^{\frac{6n}{n+1}+n-1}}{(1+r^{2})^{\frac{n^{2}}{n+1}}}\ dr \Big)^{\frac{n+1}{2n}}\notag\\
&\leq \left\{\begin{array}{ll}
\varepsilon^{\frac{n-1}{2}}, & 1\leq n\leq 6\\
\\
\varepsilon^{3}|\ln(\varepsilon)|^{\frac{4}{7}}, & n=7\\
\\
\varepsilon^{3}, & n\geq 8
\end{array}
\right..\notag
\end{align}
We finish now by estimating $B_{2}$: 
\begin{align}
\|B_{2}\|_{H^{-\frac{1}{2}}}&\leq C \Big(\int_{|x|\leq 2\delta}|x|^{\frac{6n}{n+1}}|\Psi_{\varepsilon}(x)|^{\frac{2n}{n+1}}\ dx \Big)^{\frac{n+1}{2n}}\notag\\
&\leq C  \left\{\begin{array}{ll}
\varepsilon^{\frac{n-1}{2}}, & 3\leq n\leq 8\\
\\
\varepsilon^{4}|\ln(\varepsilon)|^{\frac{5}{9}}, & n=9\\
\\
\varepsilon^{4}, & n\geq 10
\end{array}
\right..\notag
\end{align}
All the previous estimates can be summarized as follows:
\begin{equation}\label{Reps}
\|\nabla J_{g}(\varphi_{\varepsilon})\|_{H^{-\frac{1}{2}}} = \| \overline{D}\varphi_{\varepsilon}-(V_{g}*|\varphi_{\varepsilon}|^{2})\varphi_{\varepsilon} \|_{H^{-\frac{1}{2}}} \leq \left\{\begin{array}{ll}
O(\varepsilon^{\frac{n-1}{2}}), & 4\leq n\leq 6\\
\\
O(\varepsilon^{3}|\ln(\varepsilon)|^{\frac{4}{7}}), & n=7\\
\\
O(\varepsilon^{3}), & n\geq 8
\end{array}
\right. 
\end{equation}
%Combining $(\ref{phieps})$ and $(\ref{Reps})$ yields the desired result.
which is the desired result.
\end{proof}

\end{document}
